\textbf{ Fichero general de prueba }

Este fichero no contiene informacion ordenada o con una finalidad.
\textit { Solo lo uso para testar nuevas features. } \newline

Es por ello que tiene espacios, no mantiene una estructura clara y tampoco
sigue un patron particular, inclusive, agregare de forma intencional malas
practicas, con la intencion de llevar al limite este parser.

\chapter{ Listas }

\begin{enumerate}
    \item Test de enumeracion
    \item Agregamos un segundo item
    \item \begin{enumerate}
            \item Lista anidada
            \item Auto cierre de etiqueta
          \end{enumerate}
\end{enumerate}

\begin{itemize}
    \item Test de lista
    \item Prueba con inline math $\int 2x^{3y}$
    \item Cierre implicito del anterior
\end{itemize}

\begin{equation}
    E = mc^2
\end{equation}

\chapter{ Comentarios }

Dentro de cada elemento se debe restringir el contenido que se puede escribir,
ie, que pasa si dentro de una \note[nota]{ Elemento flotante que sirve para
agregar aclaraciones, referencias o incluso enlaces } intentasemos agregar % test comment
elementos como \note[equations]{ Son elementos demaciado grandes que haria que
una nota fuese muy grande}, por ello, inline_only sirve para restringir esto
%\note{test}{ Intentando romper notas con $$maths$$ no simples }

\chapter{ Imagenes }

\begin{figure}
  \includegraphics[80]{img/topologia.png}
  \caption{Prueba de imagen}
\end{figure}

\chapter{ Bloques }

Creamos ahora un bloque, que permite agrupar elementos y mostrarlos mediante botones:
\begin{block}
   \item Son utiles para mostrar distintas demostraciones del mismo elemento.
        Por ejemplo, en Topologia se puede abordar desde 2 perspectirvas, por un lado
        usar metricas y por otro lado, el concepto de bolas.
        Creo que puede ser util no tratarlos de forma indistinta puesto que, aun que
        las \note[bolas]{Una bola es un subconjunto de $R^n$} aparecen como una
        consecuencia de tener una metrica, las bolas
        por si mismas tambien son un concepto que no necesariamente requeiren de
        tener definida una metrica.
    \item Uno de los problemas al crear este entorno es la clara separacion de los
        distintos bloques.
        
        En primera instancia, se intento con un comando tipo \backslash block\{c1\}\{c2\}\{...
        Pero esto es un problema porque la cantidad de argumentos o "subbloques"
        es variable y eso hace que crear un registro se vuelva mas complejo y determinar
        el cierre tambien.
    \item Si bien, no me agrada del todo usar items, es verdad que es mucho mas simple
        usar items e identaciones para organizar el codigo
\end{block}

Y los bloques a su vez, puede contener cualquier elemento que se requiera, tanto ecuaciones
como imagenes, notas, etc.

Y ahora integramos referencias. Esto se logro con un diccionario label:\{id, index\}, asi cuando
se crea una nueva referencia, entonces tenemos la key y automaticamente, un contador. Los
contadores formales aun no se agregan.

\chapter{ Referencias }

\begin{equation}
    E = mc^2
\end{equation}
\label{eq:einstein}

Y se citan con \backslash ref\{id\}, ej, \ref{eq:einstein}

Como se agregan al output como hermano del arbol, este ya no regresa un array sino un diccionario.

Por fin se agregaron los custom components de la libreria original. Sin embargo, crear una
retrocompatibilidad implico un cambio grande en dichas librerias.

\chapter{ Customs components }

En este punto, el core ya parsea la base de latex, se pueden integrar equaciones, listas, imagenes, referencias
y algun elemento mas. Si bien, con esto ya se pueden crear documentos, tambien es cierto que la intencion de este
"parser" es expandir LaTeX de una forma sencilla. Sin tener que aprender otra sintaxis y semantica, unicamente se
agregan algunos comandos o componentes extras. Estos siguen la sintaxis de latex y como ya podemos leer y
procesar la sintaxis de latex, la adaptacion del escritor a nuevos componentes es relativamente sencilla.

\begin{axiom}
  Los custom components dejaron de ser commands inline para ser ahora
  environments. Usando key value props.
\end{axiom}

\begin{axiom}[title=Arquimedes, label=ax-arq]
  Para todo $x \in \N$ existe $y \in \N$ tal que $y > x$.
\end{axiom}

\begin{theorem}[Bolzano-Weierstrass]
  Toda sucesion acotada en $\R^n$ tiene una subsucesion convergente.
\end{theorem}

\begin{definition}[title=Espacio metrico, label=def-met]
  Un espacio metrico es un par $(X, d)$ donde ...
\end{definition}

\begin{proof}
  Sea $\epsilon > 0$ ...
  \qed
\end{proof}

\begin{exercise}[label=ej-1]
  Demuestre que ...
\end{exercise}

\section{ Capitulos }

En este punto, se integran las secciones y capitulos. Algo esencial al momento de escribir
documentos largos (libros, notas). Esto permite separar el contenido en agrupaciones con
informacion relacionada. A su vez, nos permite crear una lista de capitulos para crear
referencias y permitir al usuario viajar rapido entre contenido.
